\chapter{Multirresolución}\label{cap:multirresolucion}

El \cref{cap:stft} terminó con un dilema: una ventana larga separa las notas graves pero emborrona
los ataques; una corta hace lo contrario, y ninguna ventana resuelve las dos cosas
(\cref{tab:resolucion}). Sin embargo, los sonidos reales casi nunca necesitan las dos resoluciones
\emph{en la misma frecuencia}. En los graves, los armónicos están muy juntos en hertz y lo que
importa es separarlos. En los agudos están lejos entre sí, y lo que importa es la precisión
temporal de los transientes (platillos, consonantes, ataques). La multirresolución usa ventanas
distintas en zonas distintas del espectro.

\section{Dos caminos}

\begin{enumerate}
  \item \textbf{Transformadas no estacionarias.} La transformada de Gabor no estacionaria (NSGT)
        permite que la ventana cambie de forma continua con la frecuencia. Con ventanas
        proporcionales a $1/f$ se obtiene una transformada de Q constante invertible, con la misma
        resolución relativa en cada octava~\cite{velasco2011}. Es la solución más elegante, pero
        cambia por completo la malla: los coeficientes ya no forman una tabla rectangular.
  \item \textbf{Bandas.} Partir el espectro en unas cuantas bandas y analizar cada una con su
        propia STFT. Es más tosco, pero conserva todo lo que ya sabemos de la STFT: marcos,
        inversa de mínimos cuadrados y máscaras.
\end{enumerate}

Fourierizer sigue el segundo camino, con una simplificación: en vez de filtrar la señal antes de
analizarla, \emph{la banda se incorpora a la máscara}.

\section{Bandas como pesos en la máscara}

Sea $\{w_b\}_{b=1}^{B}$ una familia de funciones de la frecuencia con
\begin{equation}\label{eq:particion}
  w_b(f)\ge0,\qquad \sum_{b=1}^{B}w_b(f)=1\quad\text{para toda } f,
\end{equation}
una \emph{partición de la unidad}. A cada banda le asociamos una malla $\mathcal G_b$, con su
operador de análisis $\mathcal A_b$. La capa multirresolución es
\begin{equation}\label{eq:multires}
  y=\sum_{b=1}^{B}\mathcal A_b^{\dagger}\bigl((M\,w_b)\odot\mathcal A_b x\bigr),
\end{equation}
donde $M$ es la misma selección de siempre, rasterizada en cada malla. Es decir: cada banda es una
resíntesis enmascarada normal, en su propia malla, con la selección multiplicada por el peso de esa
banda.

En Fourierizer hay tres bandas. La malla de cada una es la de la capa escalada por $4$, $1$ y
$\tfrac14$, con la misma ventana, el mismo cociente $N/L$ y el mismo traslape $L/H$:
\begin{center}\small
\begin{tabular}{lcc}
\toprule
banda & frecuencias & ventana (con $N=4096$ en la capa) \\
\midrule
graves & por debajo de $\sim300\Hz$ & $16384$ ($341$ ms) \\
medios & entre $\sim300\Hz$ y $\sim3\,\mathrm{kHz}$ & $4096$ ($85$ ms) \\
agudos & por encima de $\sim3\,\mathrm{kHz}$ & $1024$ ($21$ ms) \\
\bottomrule
\end{tabular}
\end{center}
Los pesos cruzan con un coseno alzado en $\log_2 f$, de una octava de ancho y centrado en $f_c=300$ y
$3000\Hz$:
\[
  r_{f_c}(f)=S\!\Bigl(\log_2\frac{f}{f_c}+\tfrac12\Bigr),\qquad
  w_{\text{graves}}=1-r_{300},\quad w_{\text{medios}}=r_{300}-r_{3000},\quad w_{\text{agudos}}=r_{3000},
\]
con $S$ la subida de coseno de la \cref{sec:fades}. Como $r_{300}\ge r_{3000}$, los tres pesos son no
negativos y suman $1$: cumplen \eqref{eq:particion}.

\section{¿Se reconstruye la señal?}

Con la selección completa ($M\equiv1$), una buena multirresolución debería devolver la señal. Para
cada banda, $T_b x=\mathcal A_b^{\dagger}(w_b\odot\mathcal A_bx)$ aplica un peso que no depende del
tiempo. Si la respuesta al impulso $h_b$ de $w_b$ (la transformada inversa del peso) es mucho más
corta que la ventana de la malla $b$, cada fragmento ventaneado se convoluciona con $h_b$ casi sin
notar los bordes de la ventana, y $T_bx\approx h_b*x$. En ese régimen,
\[
  \sum_b T_b x\approx\Bigl(\sum_b h_b\Bigr)*x=\delta*x=x,
\]
porque la suma de los pesos es $1$ y su transformada inversa es el impulso $\delta$.

La condición ``$h_b$ corta'' equivale a que el peso varíe despacio en comparación con la resolución
en frecuencia de la malla, que es el mismo criterio de consistencia del \cref{cap:stft}. Las
transiciones de una octava se diseñaron para cumplirlo con holgura:
\begin{center}\small
\begin{tabular}{lccc}
\toprule
cruce & ancho de la transición & malla & ancho a $-6\dB$ (Hann) \\
\midrule
$300\Hz$  & $212$--$424\Hz$ ($212\Hz$)          & graves, $N=16384$ & $5.9\Hz$ \\
          &                                      & medios, $N=4096$  & $23.4\Hz$ \\
$3\,\mathrm{kHz}$ & $2.1$--$4.2\,\mathrm{kHz}$ ($2.1\,\mathrm{kHz}$) & medios, $N=4096$ & $23.4\Hz$ \\
          &                                      & agudos, $N=1024$  & $93.8\Hz$ \\
\bottomrule
\end{tabular}
\end{center}
En todos los casos, la transición es al menos veinte veces más ancha que el lóbulo principal. En
la prueba del motor de referencia (ruido y dos tonos, Hann con $N=1024$), las tres bandas
reconstruyen la entrada con un error máximo menor al 2\,\% de su pico. Los tres motores de
Fourierizer (Python, TypeScript y C++) coinciden entre sí al nivel del redondeo.

\begin{observacion}
La multirresolución aplica a la operación \emph{Sum}, donde cada capa se reconstruye por separado.
Las operaciones \emph{Max}, \emph{Cross} y \emph{Product} (\cref{cap:operaciones}) comparan
coeficientes de capas distintas celda por celda, así que necesitan una sola malla común.
\end{observacion}

\section{Cómo se dibuja}

En el espectrograma de una capa multirresolución, cada renglón sale de la imagen de la banda cuyo
peso domina esa frecuencia: los graves con la ventana larga y los agudos con la corta. Para que las
tres imágenes se puedan pegar, los niveles se comparan en dBFS y no en la escala propia de cada
imagen. Una parcial de amplitud $A$ que cae en el centro de un bin produce $|\chi|=A\sum_nw[n]/2$
(\cref{cor:parcial}), que crece con el largo de la ventana. Dividir entre $\sum_nw[n]/2$ quita esa
dependencia: la misma parcial mide lo mismo en las tres bandas. Las bandas largas se dibujan con
el salto de la capa, de modo que el emborronamiento temporal que se ve es el de la ventana larga y
no el de un muestreo grueso.

\begin{estudio}
\textbf{Cuándo usar \emph{Resolution: Multi}.}
\begin{itemize}
  \item \emph{Baterías completas, loops y mezclas.} El bombo y el bajo quedan separados en los
        graves y los platillos conservan su ataque. Es el caso típico.
  \item \emph{Voces.} Las consonantes (agudas y breves) quedan nítidas, sin perder la definición
        de los armónicos de la vocal.
  \item \emph{Pads y notas sostenidas aisladas.} Casi no hace falta: una sola ventana larga basta.
\end{itemize}
Cuesta unas tres veces más cálculo que una sola ventana, porque son tres resíntesis. Los cruces
están fijos en $300\Hz$ y $3\,\mathrm{kHz}$. Si seleccionas justo sobre un cruce, recuerda que ahí
suenan dos bandas a la vez, cada una con su resolución.
\end{estudio}
